\subsection{Đề 2}
\Opensolutionfile{ans}[ans/ans-2-B14]
\caulc
%%%=============EX_1=============%%%
\begin{ex}%[2H5N1-2]
Trong không gian $Oxyz$, cho mặt phẳng $(P)$ có phương trình $5x+3y-z+1=0$. Mặt phẳng $(P)$ có một vectơ pháp tuyến là
\choice
{$\vec{n}=(5 ; 3 ; 1)$}
{$\vec{n}=(-5 ; 3 ; 1)$}
{\True $\vec{n}=(5 ; 3 ;-1)$}
{$\vec{n}=(5 ;-3 ;-1)$}
\loigiai{Mặt phẳng $(P)$ có vectơ pháp tuyến là $\vec{n}=(5 ; 3 ;-1)$.}
\end{ex}

%%%=============EX_2=============%%%
\begin{ex}%[2H5N1-2]
Trong không gian $Oxyz$, cho ba điểm $A(1 ;-2 ; 1)$, $B(-1 ; 3 ; 3)$, $C(2 ;-4 ; 2)$. Một vectơ pháp tuyến $\vec{n}$ của mặt phẳng $(A B C)$ là
\choice
{\True $\vec{n}=(9 ; 4 ;-1)$}
{$\vec{n}=(9 ; 4 ; 1)$}
{$\vec{n}=(4 ; 9 ;-1)$}
{$\vec{n}=(-1 ; 9 ; 4)$}
\loigiai{
Ta có $\vec{A B}=(-2 ; 5 ; 2)$, $\vec{A C}=(1 ;-2 ; 1)$
$\Leftrightarrow \vec{n}_{(ABC)}=\left[\vec{A B}, \vec{A C} \right]=(9 ; 4 ;-1)$.
}
\end{ex}

%%%=============EX_3=============%%%
\begin{ex}%[2H5N1-3]
Trong không gian $Oxyz$, phương trình mặt phẳng $(P)$ đi qua điểm $A(-1 ; 2 ; 0)$ và nhận $\vec{n}=(-1 ; 0 ; 2)$ là vectơ pháp tuyến có phương trình là
\choice
{$-x+2 y-5=0$}
{$-x+2 z-5=0$}
{$-x+2 y-5=0$}
{\True $-x+2 z-1=0$}
\loigiai{
Mặt phẳng $(P)$ đi qua điểm $A(-1 ; 2 ; 0)$ và nhận $\vec{n}=(-1 ; 0 ; 2)$ là vectơ pháp tuyến có phương trình là
$$-1(x+1)+0(y-2)+2(z-0)=0 \Leftrightarrow-x-1+2 z=0 \Leftrightarrow-x+2 z-1=0.$$
Vậy $(P)\colon-x+2 z-1=0$.
}
\end{ex}

%%%=============EX_4=============%%%
\begin{ex}%[2H5H1-3]
Trong không gian $Oxyz$, cho hai điểm $A(-1 ; 0 ; 1)$, $B(-2 ; 1 ; 1)$. Phương trình mặt phẳng trung trực của đoạn $A B$ là
\choice
{$x-y-2=0$}
{$x-y+1=0$}
{\True $x-y+2=0$}
{$-x+y+2=0$}
\loigiai{
Ta có $\vec{A B}=(-1 ; 1 ; 0)$. Trung điểm $I$ của đoạn $A B$ là $I\left(-\dfrac{3}{2} ; \dfrac{1}{2} ; 1\right)$.\\
Mặt phẳng trung trực của đọan $\mathrm{AB}$ là $-\left(x+\dfrac{3}{2}\right)+\left(y-\dfrac{1}{2}\right)=0$ hay $x-y+2=0$.
}
\end{ex}

%%%=============EX_5=============%%%
\begin{ex}%[2H5H1-4]
Trong không gian $Oxyz$, cho hai mặt phẳng $(\alpha)\colon  3 x+(m-1) y+4 z-2=0$, $ (\beta)\colon  n x+(m+2) y+2 z+4=0$. Với giá trị thực của $m$, $n$ bằng bao nhiêu để $(\alpha)$ song song $(\beta)$.
\choice
{$m=3$; $n=-6$}
{$m=3$; $n=6$}
{$m=-3$; $n=-6$}
{\True $m=-3$; $n=6$}
\loigiai{
Đề $(\alpha)$ song song $(\beta) \Leftrightarrow \dfrac{3}{n}=\dfrac{m-1}{m+2}=\dfrac{4}{2} \neq \dfrac{4}{-2} \Leftrightarrow m=-3 ; \, n=6$.\\
Vậy $m=-3$; $n=6$.
}
\end{ex}

%%%=============EX_6=============%%%
\begin{ex}%[2H5H1-3]
Trong không gian $O x y z$, phương trình mặt phẳng $(P)$ đi qua điểm $M(-2 ; 3 ; 1)$ và song song với mặt phẳng $(Q)\colon  x+3 y-2 z+2=0$ là
\choice
{$(Q)\colon  x+3 y-2 z+5=0$}
{\True $(Q)\colon  x+3 y-2 z-5=0$}
{$(Q)\colon  -2x+3 y+z+2=0$}
{$(Q)\colon  -2x+3 y+z-5=0$}
\loigiai{
Mặt phẳng $(P)$ song song với mặt phẳng $(Q)\colon  x+3 y-2 z+2=0$ nên mặt phẳng $(P)$ có phương trình dạng $(Q)\colon  x+3 y-2 z+m=0\ (m\neq2)$.\\
Mặt phẳng $(P)$ đi qua điểm $M(-2 ; 3 ; 1)$ nên $-2+3\cdot3-2\cdot1+m=0 \Leftrightarrow m=-5$ (thỏa mãn).\\
Vậy phương trình mặt phẳng $(P)$ là $(Q)\colon  x+3 y-2 z-5=0$.
}
\end{ex}

%%%=============EX_7=============%%%
\begin{ex}%[2H5H1-3]
Trong không gian $Oxyz$, cho ba điểm $A(3 ;-2 ;-2)$, $B(3 ; 2 ; 0)$, $C(0 ; 2 ; 1)$. Phương trình mặt phẳng $(A B C)$ là
\choice
{$4 y+2 z-3=0$}
{$3 x+2 y+1=0$}
{\True $2 x-3 y+6 z=0$}
{$2 y+z-3=0$}
\loigiai{
Ta có $\vec{A B}=(0 ; 4 ; 2)$, $\vec{A C}=(-3 ; 4 ; 3)$, $\left[\vec{A B},\vec{A C}\right]=(4 ;-6 ; 12)$.\\
$(A B C)$ qua $A(3 ;-2 ;-2)$ và có vectơ pháp tuyến $\vec{n}=(2 ;-3 ; 6)$.\\
$(A B C)\colon  2(x-3)-3(y+2)+6(z+2)=0 $ hay $(A B C)\colon 2 x-3 y+6 z=0.$
}
\end{ex}

%%%=============EX_8=============%%%
\begin{ex}%[2H5H1-3]
Trong không gian $Oxyz$, cho mặt phẳng $(\alpha)$ đi qua $A(2 ;-1 ; 4), B(3 ; 2 ;-1)$ và vuông góc với mặt phẳng $(Q)\colon  x+y+2 z-3=0$. Phương trình mặt phẳng $(\alpha)$ là
\choice
{\True $5x+3y-4z+9=0$}
{$x+3y-5z+21=0$}
{$x+y+2z-3=0$}
{$5x+3y-4z=0$}
\loigiai{
Ta có $\vec{A B}=(1 ; 3 ;-5)$, $\vec{n}_{(Q)}=(1 ; 1 ; 2)$, $\left[\vec{A B}, \vec{n}_{(Q)}\right]=(-10 ;-6 ; 8)$.\\
Mặt phẳng $(\alpha)$ đi qua $A(2 ;-1 ; 4)$ và có vectơ pháp tuyến 
$ \vec{n}=(5 ; 3 ;-4)$ có phương trình là\\
$(\alpha)\colon 5(x-2)+3(y+1)-4(z-4)=0 $ hay $(\alpha)\colon  5x+3y-4z+9=0$.
}
\end{ex}

%%%=============EX_9=============%%%
\begin{ex}%[2H5H1-3]
Trong không gian $Oxyz$, mặt phẳng $(\alpha)$ đi qua điểm $A(2 ;-1 ; 5)$ và vuông góc với hai mặt phẳng $(P)\colon  3 x-2 y+z+7=0$ và $(Q)\colon  5 x-4 y+3 z+1=0$. Phương trình mặt phẳng $(\alpha)$ là
\choice
{\True $x+2 y+z-5=0$}
{$2 x-4 y-2 z-10=0$}
{$2 x+4 y+2 z+10=0$}
{$x+2 y-z+5=0$}
\loigiai{
Mặt phằng $(P)$ có một vectơ pháp tuyến là $\vec{n}_{(P)}=(3 ;-2 ; 1)$.\\
Mặt phẳng $(Q)$ có một vectơ pháp tuyến là $\vec{n}_{(Q)}=(5 ;-4 ; 3)$.\\
Mặt phẳng $(\alpha)$ vuông góc với 2 mặt phẳng $(P)\colon  3 x-2 y+z+7=0,(Q)\colon  5 x-4 y+3 z+1=0$ nên có một vectơ pháp tuyến là $\vec{n}_{(\alpha)}=\left[\vec{n}_{(P)}, \vec{n}_{(Q)}\right]=(-2 ;-4 ;-2)$.\\
Vậy phương trình mặt phẳng $(\alpha)$ là $x+2 y+z-5=0$.
}
\end{ex}

%%%=============EX_10=============%%%
\begin{ex}%[2H5H1-5]
Trong không gian $Oxyz$, tọa độ điểm $M$ nằm trên trục $O y$ và cách đều hai mặt phẳng $(P)\colon  x+y-z+1=0$ và $(Q)\colon  x-y+z-5=0$ là
\choice
{\True $M(0 ;-3 ; 0)$}
{$M(0 ; 3 ; 0)$}
{$M(0 ;-2 ; 0)$}
{$M(0 ; 2 ; 0)$}
\loigiai{
Ta có $M \in Oy $ nên $ M(0 ; m ; 0)$.\\
Giả thiết có $\mathrm{d}(M,(P))=\mathrm{d}(M,(Q)) \Leftrightarrow \dfrac{|m+1|}{\sqrt{3}}=\dfrac{|-m-5|}{\sqrt{3}} \Leftrightarrow m=-3$.\\
Vậy $M(0 ;-3 ; 0)$.
}
\end{ex}

%%%=============EX_11=============%%%
\begin{ex}%[2H5H1-5]
Trong không gian $O x y z$, phương trình mặt phẳng $(P)$ song song với mặt phẳng $(Q)\colon  x+2 y-2 z+1=0$ và cách $(Q)$ một khoảng bằng $3$ là
\choice
{$(P)\colon  x+2 y-2 z-8=0$}
{$(P)\colon  x+2 y-2 z-10=0$}
{$(P)\colon  x+2 y-2 z+8=0$, $(P)\colon  x+2 y-2 z-10=0$}
{\True $(P)\colon  x+2 y-2 z-8=0$, $(P)\colon  x+2 y-2 z+10=0$}
\loigiai{
Trên mặt phẳng $(Q)\colon  x+2 y-2 z+1=0$ chọn điểm $M(-1 ; 0 ; 0)$.\\
Do $(P)\parallel (Q)$ nên phương trình của mặt phẳng $(P)$ có dạng $x+2 y-2 z+D=0$ với $D \neq 1$.\\
Vì $\mathrm{d}((P),(Q))=3 \Leftrightarrow \mathrm{d}(M,(P))=3 \Leftrightarrow \dfrac{|-1+D|}{\sqrt{1^{2}+2^{2}+(-2)^{2}}}=3 \Leftrightarrow|-1+D|=9 \Leftrightarrow \hoac{&D=-8\\&D=10.}$\\
Vậy $(P)\colon  x+2 y-2 z-8=0$, $(P)\colon  x+2 y-2 z+10=0$.
}
\end{ex}

%%%=============EX_12=============%%%
\begin{ex}%[2H5H1-3]
Trong không gian $Oxyz$, cho hình lăng trụ $ABC.A'B'C'$ với $A(5 ; 1 ; 3)$, $B(1 ; 2 ; 6)$, $C(5 ; 0 ; 4)$, $A'(4 ; 0 ; 6)$. Viết phương trình mặt phẳng $(A'B'C')$.
\choice
{$x+y+z-9=0$}
{\True $x+y+z-10=0$}
{$x+y+z-8=0$}
{$x+2 y+z-10=0$}
\loigiai{
Mặt phẳng $(A'B'C')$ nhận $\vec{AB}=(-4;1;3)$, $\vec{AC}=(0;-1;1)$ làm cặp vectơ chỉ phương nên có vectơ pháp tuyến là $\vec n=\left[ \vec{AB},\vec{AC}\right] =(4;4;4).$\\
Mặt phẳng $(A'B'C')$ đi qua $A'(4;0;6)$ và nhận $\vec n=(1;1;1)$ làm một vectơ pháp tuyến nên có phương trình
$x+y+z-10=0$.
}
\end{ex}

\Closesolutionfile{ans}
\indapan{10}{ans/ans-2-B14}

\cauds
\Opensolutionfile{ans}[ans/ans-2-B14-DS]

%%%=============EX_13=============%%%
\begin{ex}%[2H5H1-3]
Trong không gian $Oxyz$, cho mặt phẳng $(P)\colon 3x+2y-4z+1=0$. 
\choiceTF
{\True $\vec{n}=(3;2;-4)$ là một vectơ pháp tuyến của $(P)$}
{\True $A(-1;-1;-1)$ là một điểm nằm trên $(P)$}
{$(P)$ song song với mặt phẳng $(Q)\colon 6x+4y+8z=0$}
{$(P)$ vuông góc với mặt phẳng $(R)\colon 2x+y-z+3=0$}
\loigiai{
\begin{itemchoice}
\itemch Đúng.
\itemch Đúng. Vì $3(-1)+2(-1)-4(-1)+1=0$.
\itemch Sai. Vì với $\vec{n}_{(P)}=(3;2;-4)$ và $\vec{n}_{(Q)}=(6;4;8)$ lần lượt là vectơ pháp tuyến của $(P)$ và $(Q)$, ta thấy $\dfrac{3}{6}=\dfrac{2}{4}\ne -\dfrac{4}{8}$ nên $\vec{n}_{(P)}$ không cùng phương với $\vec{n}_{(Q)}$. Do đó $(P)$ không song song với $(Q)$.
\itemch Sai. Vì với $\vec{n}_{(P)}=(3;2;-4)$ và $\vec{n}_{(R)}=(2;1;-1)$ lần lượt là vectơ pháp tuyến của $(P)$ và $(R)$, ta thấy $\vec{n}_{(P)}\cdot\vec{n}_{(R)}=12\neq 0$. Do đó $(P)$ không vuông góc với $(R)$.
\end{itemchoice}
}
\end{ex}

%%%=============EX_14=============%%%
\begin{ex}%[2H5H1-4]
Trong không gian $Oxyz$, cho mặt phẳng $(P)$ đi qua $A(3;-2;5)$ và có vectơ pháp tuyến $\vec{n}=(4;-3;2)$. 
\choiceTF
{Phương trình mặt phẳng $(P)$ là $3x-2y+5z-28=0$}
{$(P)$ song song với mặt phẳng $(Q)\colon 2x+11y+8z-5=0$}
{\True $(P)$ vuông góc với mặt phẳng $(R)\colon x+2y-z+3=0$}
{Khoảng cách từ điểm $B(1;1;-1)$ đến mặt phẳng $(P)$ là $5$}
\loigiai{
\begin{itemchoice}
\itemch Sai. Phương trình mặt phẳng $(P)$ đi qua $A(3;-2;5)$ có một vectơ pháp tuyến $\vec{n}_{(P)}=(4;-3;2)$ có dạng $4(x-3)-3(y+2)+2(z-5)=0$, hay $(P)\colon 4x-3y+2z-28=0$.
\itemch Sai. Vì với $\vec{n}_{(P)}=(4;-3;2)$ và $\vec{n}_{(Q)}=(2;11;8)$ lần lượt là vectơ pháp tuyến của $(P)$ và $(Q)$, ta thấy $\dfrac{4}{2}\ne-\dfrac{3}{11}$ nên nên $\vec{n}_{(P)}$ không cùng phương với $\vec{n}_{(Q)}$. Do đó $(P)$ không song song với $(Q)$.
\itemch Đúng. Vì với $\vec{n}_{(P)}=(4;-3;2)$ và $\vec{n}_{(R)}=(1;2;1)$ lần lượt là vectơ pháp tuyến của $(P)$ và $(R)$, ta thấy $\vec{n}\cdot \vec{n}_2=0$ nên $(P)\perp (R)$.
\itemch Sai. Ta có
$\mathrm{d}\left(B,(P)\right)=\dfrac{\left|4\cdot 1-3\cdot 1+2\cdot (-1)-28\right|}{\sqrt{4^2+(-3)^2+2^2}}=\sqrt{29}.$
\end{itemchoice}
}
\end{ex}

%%%=============EX_15=============%%%
\begin{ex}%[2H5H1-5]
Trong không gian $Oxyz$, cho tứ diện $ ABCD $ có tọa độ các đỉnh $ A(5;1;3)$, $B(1;6;2)$, $C(5;0;4)$, $D(4;0;6) $.
\choiceTF
{\True Phương trình mặt phẳng $(ACD)$ là $2x+y+z-14=0$}
{Khoảng cách từ điểm $B$ đến mặt phẳng $(ACD)$ là $\dfrac{2\sqrt{3}}{3}$}
{Phương trình mặt phẳng $(P)$ đi qua điểm $A$ và song song với $(BCD)$ là $3x+2y+17=0$}
{\True Phương trình mặt phẳng $ (Q) $ chứa cạnh $ AB $, song song với cạnh $ CD $ là $10x+9y+5z-74=0$}
\loigiai{
Ta có $ \vec{AC}=(0;-1;1)$, $\vec{AD}=(-1;-1;3)$, $\vec{BC}=(4;-6;2)$, $\vec{BD}=(4;-6;4) $.
\begin{itemchoice}
\itemch Đúng. Mặt phẳng $ (ACD) $ qua $ A(5;1;3) $ và có vectơ pháp tuyến $ \vec{n}=\left[\vec{AC},\vec{AD}\right]=(-2;-1;-1) $ có phương trình là 
$ -2(x-5)-(y-1)-(z-3)=0\Leftrightarrow 2x+y+z-14=0 .$
\itemch Sai. Ta có $\mathrm{d}(B,(ACD)=\dfrac{|2\cdot 1+6+2-14|}{\sqrt{2^2+1^2+1^2}}=\dfrac{2\sqrt{6}}{3}$.
\itemch Sai. Mặt phẳng $(P)\parallel (BCD)$ nên có vectơ pháp tuyến $ \vec{n}=\left[\vec{BC},\vec{BD}\right]=(-12;-8;0) $.\\
Khi đó phương trình mặt phẳng $(P)$ có dạng $3x+2y+c=0$.\\
Do $A\in (P)$ nên $3\cdot5+2\cdot1+c=0\Leftrightarrow c=-17$. Hay $(P)\colon 3x+2y-17=0$.
\itemch Đúng. Ta có $ \vec{AB}=(-4;5;-1) $, $\vec{CD}=(-1;0;2)$.\\
Mặt phẳng $ (\alpha) $ chứa cạnh $ AB $ và song song với cạnh $ CD $ qua $ A(5;1;3) $ và có vectơ pháp tuyến $ \vec{n}=\left[\vec{AB},\vec{CD}\right]=(10;9;5) $ có phương trình là 
$$ 10(x-5)+9(y-1)+5(z-3)=0\Leftrightarrow 10x+9y+5z-74=0 .$$
\end{itemchoice} 
}
\end{ex}
%%%=============EX_16=============%%%
\begin{ex}%[2H5V1-5]
Trong không gian $Oxyz$, cho hai điểm $A(1;3;-4)$ và $B(-1;2;2)$. 
\choiceTF
{Phương trình mặt phẳng trung trực của đoạn $AB$ là $-2x+y+6z+\dfrac{17}{2}=0$}
{\True Tọa độ hình chiếu vuông góc của điểm $M\left(5;\dfrac{5}{2};-13\right)$ lên mặt phẳng trung trực của đoạn $AB$ là $H\left(1;\dfrac{1}{2};-1\right)$}
{Phương trình mặt phẳng $(Q)$ đi qua hai điểm $A$, $B$ và vuông góc với mặt phẳng $(R)\colon x-3y+2z+4=0$ là $4x-2y-5z+9=0$}
{\True Phương trình mặt phẳng $(\alpha)$ đi qua hai điểm $A$, $B$ sao cho khoảng cách từ $C(2;11;-4)$ đến $(\alpha)$ bằng $6$ là $(\alpha)\colon x+y+\dfrac{1}{2}z-2=0$ hoặc $(\alpha)\colon -9x+13y-\dfrac{5}{6}z-\dfrac{100}{3}=0$}
\loigiai{
\begin{itemchoice}
\itemch Sai. Ta có $M\left(0;\dfrac{5}{2};-1\right)$ là trung điểm của $AB$. Khi đó mặt phẳng trung trực $(P)$ của $AB$ đi qua $M\left(0;\dfrac{5}{2};-1\right)$, có vectơ pháp tuyến $\vec{AB}=(-2;-1;6)$ có phương trình là
$$-2(x-0)-1\left(y-\dfrac{5}{2}\right)+6(z+1)=0\Leftrightarrow -2x-y+6z+\dfrac{17}{2}=0.$$
\itemch Đúng. Gọi $H(x;y;z)$ là hình chiếu cần tìm, ta có $\vec{MH}=\left(x-5;y-\dfrac{5}{2};z+13\right)$.\\
Ta có $(P)\colon -2x-y+6z+\dfrac{17}{2}=0$ là mặt phẳng trung trực của $AB$.\\
Ta có $\vec{MH}\perp (P)$ nên $\vec{MH}$ cùng phương với vectơ pháp tuyến $\vec{n}=(-2;-1;6)$ của $(P)$, hơn nữa $H\in (P)$. Do đó tọa độ điểm $H$ là nghiệm của hệ phương trình
$$\heva{&\dfrac{x-5}{-2}=\dfrac{5}{2}-y\\&\dfrac{5}{2}-y=\dfrac{z+13}{6}\\&-2x-y+6z+\dfrac{17}{2}=0}
\Leftrightarrow \heva{&x-2y=0\\&-6y-z=-2\\&-2x-y+6z=-\dfrac{17}{2}} \Leftrightarrow \heva{&x=1\\&y=\dfrac{1}{2}\\&z=-1.}$$
Vậy $H\left(1;\dfrac{1}{2};-1\right)$.
\itemch Sai. Ta có $\vec{n}_{(R)}=(1;-3;2)$ là vectơ pháp tuyến của $(R)$, khi đó $\vec{n}_{(Q)}=\left[\vec{n}_{(R)},\vec{AB}\right]=(-16;-10;-7)$.\\
Mặt phẳng $(Q)$ đi qua $A(1;3;-4)$ và có vectơ pháp tuyến $\vec{n}_{(Q)}=(-16;-10;-7)$ có phương trình là
$$-16(x-1)-10\left(y-3\right)-7(z+4)=0
\Leftrightarrow -16x-10y-7z+18=0.$$
\itemch Đúng. Phương trình mặt phẳng $(\alpha)$ đi qua $A(1;3;-4)$ và có vectơ pháp tuyến $\vec{n}_{(\alpha)}=(A;B;C)$ có dạng $A(x-1)+B(y-3)+C(z+4)=0\Leftrightarrow Ax+By+Cz-A-3B+4C=0.$\\
Vì $B(-1;2;2)\in (\alpha)$ nên $-2A-B+6C=0$. (1)\\
Mặt khác $\mathrm{d}\left(C,(\alpha)\right)=6\Leftrightarrow \dfrac{\left|A+8B\right|}{\sqrt{A^2+B^2+C^2}}=6\Leftrightarrow 35A^2-16AB-28B^2+36C^2=0$. (2)\\
Từ (1) và (2), ta có
$39A^2-12AB-27B^2=0\Leftrightarrow \hoac{&A=B\\&A=-\dfrac{9}{13}B.}$\\
Với $A=B$, ta chọn $A=1$, khi đó $A=B=1,C=\dfrac{2A+B}{6}=\dfrac{1}{2}$. Ta có
$(\alpha)\colon x+y+\dfrac{1}{2}z-2=0.$\\
Với $A=-\dfrac{9}{13}B$, ta chọn $A=-9$, khi đó $B=13,C=\dfrac{2A+B}{6}=-\dfrac{5}{6}$. Ta có
\begin{eqnarray*}
(\alpha)\colon -9x+13y-\dfrac{5}{6}z-\dfrac{100}{3}=0.
\end{eqnarray*}
\end{itemchoice}
}
\end{ex}
\Closesolutionfile{ans}
\indapan{3}{ans/ans-2-B14-DS}

\Opensolutionfile{ans}[ans/ans-2-B14-KQ]
\caukq

%%%=============EX_17=============%%%
\begin{ex}%[2H5H1-3]
Trong không gian $Oxyz$, phương trình mặt phẳng $(P)$ đi qua hai điểm $A(1;2;-2)$, $B(2;4;1)$ và vuông góc với mặt phẳng $(Q)\colon x+3y+z-1=0$ có dạng $ax+by+cz-5=0$. Giá trị $a+b+c$ bằng bao nhiêu?
\shortans{$4$}
\loigiai{
Mặt phẳng $(Q)$ có vectơ pháp tuyến $\vec{n}_{(Q)}=(1;3;1)$. Mặt phẳng $(P)$ đi qua $A$, $B$ và vuông góc với $(Q)$ nên có cặp vectơ chỉ phương là $\vec{AB}=(1;2;3)$ và $\vec{n}_{(Q)}=(1;3;1)$. Do đó $(P)$ có vectơ pháp tuyến là $\vec{n}_{(P)}=\left[\vec{AB},\vec{n}_{(Q)}\right]=(-7;2;1)$.\\
Mặt phẳng $(P)$ đi qua $A(1;2;-2)$ và có vectơ pháp tuyến $\vec{n_P}=(-7;2;1)$ nên có phương trình
$$-7(x-1)+2(y-2)+1(z+2)=0\Leftrightarrow 7x-2y-z-5=0.$$
Vậy $a=7$, $b=-2$, $c=-1$. Do đó $a+b+c=4$.
}
\end{ex}

%%%=============EX_18=============%%%
\begin{ex}%[2H5H1-3]
Trong không gian $Oxyz$, mặt phẳng $(P)$ đi qua điểm $M(1;2;3)$ và cắt trục tọa độ tại $A$, $B$, $C$ (khác gốc tọa độ) sao cho $M$ là trọng tâm tam giác $ABC$ có phương trình là $ax+by+cz-18=0$. Giá trị của $abc$ bằng bao nhiêu?
\shortans{$36$}
\loigiai
{
Gọi $A(p;0;0)$, $B(0;q;0)$, $C(0;0;r)$.\\
Vì $M(1;2;3)$ là trọng tâm $\triangle ABC$ nên
$$\heva{&x_G=\dfrac{x_A+x_B+x_C}{3}\\ &y_G=\dfrac{y_A+y_B+y_C}{3}\\ &z_G = \dfrac{z_A+z_B+z_C}{3}} \Leftrightarrow \heva{&1=\dfrac{p}{3}\\ &2=\dfrac{q}{3}\\ &3=\dfrac{r}{3}} \Leftrightarrow \heva{&a=3\\ &b=6\\ &c=9.}$$
Mặt phẳng $(P)$ qua ba điểm $A(3;0;0)$, $B(0;6;0)$, $C(0;0;9)$ có dạng $\dfrac{x}{3} + \dfrac{y}{6} + \dfrac{z}{9} = 1$. \\
Suy ra $(P)\colon 6x+3y+2z-18=0$.\\
Do đó $a=6$, $b=3$ và $z=2$ nên $abc=36$.
}
\end{ex}

%%%=============EX_19=============%%%
\begin{ex}%[2H5V1-5]
Trong không gian $Oxyz$, cho hình hộp $ ABCD.A'B'C'D' $có $ A(1;0;1)$, $B(2;1;2)$, $D(1;-1;1)$, $C'(4;5;-5) $. Chiều cao của hình hộp $ ABCD.A'B'C'D' $ là $\dfrac{a\sqrt{b}}{2}$. Giá trị của $a\cdot b$ bằng bao nhiêu?
\shortans{$18$}
\loigiai{
Ta có$\vec{AB}=(1;1;1)$, $\vec{AD}=(0;-1;0) $ nên $ \vec{n}_{(ABCD)}=\left[\vec{AB},\vec{AD}\right]=(1;0;-1) $ là một vectơ pháp tuyến của $ (ABCD) $.\\ 
Mà mặt phẳng $ (ABCD) $ qua $ A(1;0;1) $ nên có phương trình là
$$ (x-1)+0(y-0)-(z-1)=0\Leftrightarrow x-z=0 .$$
Mặt phẳng $ (A'B'C'D') \parallel (ABCD)$ nên phương trình có dạng $ x-z+m=0\,\,(m\ne0) $.\\
Mà $ C'(4;5;-5) \in(A'B'C'D')\Leftrightarrow 1\cdot4+(-1)\cdot(-5)+m=0\Leftrightarrow m=-9$ (nhận).\\
Vậy phương trình mặt phẳng $ (A'B'C'D') $ là $ x-z-9=0 $.\\
Gọi $ h $ là chiều cao của hình hộp, suy ra $  h=\mathrm{d}(A,(A'B'C'D')) =\dfrac{|1\cdot1+0-1\cdot1-9|}{\sqrt{1^2+0^2+1^2}}=\dfrac{9\sqrt{2}}{2}$.\\
Vậy $a=9$, $b=2$. Do đó $a\cdot b=18$.
}
\end{ex}

%%%=============EX_20=============%%%
\begin{ex}%[2H5V1-4]
Trong không gian $Oxyz$, cho $(\alpha) \colon y+2z-4=0$, $(\beta) \colon  x + y -z +3 =0$, $(\gamma) \colon x+y+z-2=0$. Phương trình mặt phẳng $(P)$ qua giao tuyến hai mặt phẳng $(\alpha)$ và $(\beta)$, đồng thời $(P)$ vuông góc với mặt phẳng $(\gamma)$ có dạng $ax+by+cz+13=0$. Giá trị $a+b+c$ bằng bao nhiêu?  
\shortans{$0$}
\loigiai{
Phương trình $(P) \colon  m(y+2z-4) + n(x+y-z+3) = 0$ với $m^2 + n^2  \ne 0$.\\
Hay $(P) \colon nx + (m+n)y + (2m-n)z -4m+3n = 0.$\\
Vectơ pháp tuyến của $(\alpha)$, $(\beta)$, $(\gamma)$ lần lượt là $\vec{n}_{(\alpha)} = (0;1;2)$, $\vec{n}_{(\beta)}=(1;1;-1)$, $\vec{n}_{(\gamma)} = (1;1;1)$.\\
Do $(P)$ qua giao tuyến của $(\alpha)$ và $(\beta)$ nên 
$\vec{a}_{1}=\left[\vec{n}_{(\alpha)},\vec{n}_{(\beta)}\right] = (-3;2;-1)$ là một vectơ chỉ phương của $(P)$.\\
Do $(P) \perp (\gamma)$ nên $\vec{a}_{2} = \vec{n}_{(\gamma)} = (1;1;1)$ là một vectơ chỉ phương của $(P)$.\\
Từ đó, $\left[\vec{a}_{1},\vec{a}_{2}\right] = (3;2;-5) = \vec{n}_{(P)}$.\\
Ta có hệ phương trình $\heva{&n = 3\\&m+n = 2\\&2m-n=-5} \Leftrightarrow \heva{&m=-1\\&n=3.}$\\
Vậy $(P) \colon 3x + 2y -5z + 13 = 0$. \\
Do đó $a=3$, $b=2$, $c=-5$. Suy ra $a+b+c=0$.}
\end{ex}


%%%=============EX_21=============%%%
\begin{ex}%[2H5C1-7]
Người ta thiết kế một mái che hình chữ nhật $ ABCD $ phía trên sân khấu. Gắn hệ trục tọa độ $ Oxyz $ (đơn vị trên trục là mét), người ta xác định được toạ dộ của các điểm như sau: $ A(0;0;8)$, $B(0;20;8)$, $D(15;0;14)$, $C(15;20;14) $. Một cổng chào hình chữ nhật $ EFHG $ với tọa độ điểm $ G(8;0;4) $ dựng vuông góc với mặt đất. Người ta muốn làm các đoạn dây nối thanh ngang $ GE $ với mái che để gắn hoa và đèn led. Độ dài ngắn nhất của mỗi đoạn dây này bằng bao nhiêu mét? (làm tròn đến chữ số thập phân thứ nhất)
\begin{center}
	\includegraphics[scale=.3]{images/2P5-1-H5-16}
\end{center}
\shortans{$1{,}8$}
\loigiai{
Ta có $ A(0;0;8)$, $B(0;20;8)$, $D(15;0;14)$, $C(15;20;14) $.\\
Ta có $ \vec{AB}=(0;20;0) $, $\vec{AC}=(15;20;6)$ nên $ \vec{n}_1=\left[\vec{AB},\vec{AC}\right]=(80;0;300) $ là vectơ pháp tuyến của $ (ABCD) $.\\
Mà mặt phẳng mái che $ (ABCD) $ qua $ A(0;0;8)$ nên có phương trình
$$ 80(x-0)+0(y-0)+300(z-8)=0\Leftrightarrow 4x+15z-120=0 .$$
Độ dài ngắn nhất của dây nối thanh ngang $ GE $ với mái che là khoảng cách từ $ G $ đến mái che (mặt phẳng $ ABCD $) là $$ \mathrm{d}(G,(ABCD))=\dfrac{|4\cdot8+0+15\cdot4-120|}{\sqrt{4^2+0^2+15^2}}=\dfrac{28}{\sqrt{241}}=1{,}8\ (\text{m}). $$
}
\end{ex}
%%%=============EX_22=============%%%
\begin{ex}%[2H5C1-7]
Hai học sinh đang chuyền bóng. Bạn nữ ném bóng cho bạn nam. Quả bóng bay trên không, lệch sang phải và rơi xuống tại vị trí cách bạn nam $3$ m, cách bạn nữ $5$ m. Cho biết quỹ đạo của quả bóng nằm trong mặt phẳng $(P)$ vuông góc với mặt đất. Phương trình của $(P)$ trong không gian $Oxyz$ được mô tả như trong hình vẽ có dạng $ax+3y+cz+d=0$. Giá trị của $a+cd$ là bao nhiêu?
\begin{center}
	\includegraphics[scale=0.6]{images/12-SGK-CTST-5-1-2} 
\end{center}
\shortans{$-4$}
\loigiai{
Gọi $A(3;4;0)$ là điểm mà quả bóng rơi.\\
Mặt phẳng $(P)$ vuông góc với mặt đất $(Oxy)$ nên $(P)$ có một vectơ chỉ phương là $\Vec{u}=(0,0,1)$.\\
Ngoài ra, $(P)$ có vec tơ chỉ phương $\Vec{OA}=(3;4;0)$.\\
vectơ pháp tuyến của $(P)$ là $ \vec{n}=\left[\vec{u},\vec{OA}\right]=(-4;3;0)$.\\
Phương trình mặt phẳng $(P)$ đi qua điểm $A(3;4;0)$ và có vectơ pháp tuyến là $ \vec{n}=(-4;3;0)$ là $-4(x-3)+3(y-4)=0 \Leftrightarrow -4x+3y=0$.\\
Do đó $a=-4$, $c=d=0$. Vậy $a+cd=-4$.
}
\end{ex}

\Closesolutionfile{ans}
\indapan{6}{ans/ans-2-B14-KQ}